\section{Results}
\label{sec:results}

\begin{table*}[t]
\centering
\caption{MovieLens NDCG@10 (LGBMRanker, mean $\pm$ 95\% CI over 5 seeds).
Columns add scorer features left to right under popularity-matched
negatives; the last column uses uniform negatives. Popularity and PureSVD
use no item text.}
\label{tab:main-results}
\begin{tabular}{lcccc}
\toprule
 & \multicolumn{3}{c}{Popularity-matched negatives} & \\
\cmidrule(lr){2-4}
Pipeline & Item only & +Affinity & +Aff.+Pop. & Uniform neg. \\
\midrule
Popularity & 0.379$\pm$0.002 & 0.379$\pm$0.002 & 0.379$\pm$0.002 & 0.731$\pm$0.006 \\
PureSVD & 0.572$\pm$0.008 & 0.572$\pm$0.008 & 0.572$\pm$0.008 & 0.757$\pm$0.005 \\
\midrule
TF-IDF & 0.383$\pm$0.004 & 0.573$\pm$0.007 & 0.576$\pm$0.008 & 0.784$\pm$0.005 \\
SBERT (22M) & 0.380$\pm$0.003 & 0.590$\pm$0.007 & 0.589$\pm$0.008 & 0.785$\pm$0.005 \\
mxbai (335M) & 0.379$\pm$0.004 & 0.586$\pm$0.005 & 0.583$\pm$0.005 & 0.785$\pm$0.005 \\
phi4-mini raw (3.8B) & 0.379$\pm$0.005 & 0.486$\pm$0.006 & 0.486$\pm$0.006 & 0.747$\pm$0.006 \\
Qwen3-Emb.\ (4B) & 0.378$\pm$0.004 & 0.593$\pm$0.007 & 0.587$\pm$0.005 & 0.785$\pm$0.005 \\
Hybrid & 0.380$\pm$0.004 & 0.579$\pm$0.007 & 0.579$\pm$0.009 & 0.783$\pm$0.005 \\
\bottomrule
\end{tabular}

\end{table*}

\subsection{An item-only scorer is a popularity model}
\label{sec:movielens-results}
The ``Item only'' column of Table~\ref{tab:main-results} is the common
shortcut of feeding item embeddings straight into a ranker: the scorer
sees item vectors and nothing about the user. Under popularity-matched
negatives every encoder lands within 0.004 NDCG@10 of plain popularity
(0.378--0.383 vs.\ 0.379), and all are 0.19 below PureSVD. A scorer with no user
input can only learn one global score per item, so it ranks every
user's candidates identically and the item representation is reduced to
a predictor of popularity. This explains why every encoder looked alike
in that setting, and a study run only in this configuration would
conclude, wrongly, that encoder choice does not matter. On Amazon and
Yelp the item-only scorer falls \emph{below} popularity
(Table~\ref{tab:amazon-results}). Once the scorer receives the user--item affinity
features (``+Affinity''), TF-IDF, SBERT, mxbai, and Qwen3 gain 0.19--0.22
(phi4-mini only 0.11), and the encoders separate. Adding the popularity features changes no pipeline
by more than 0.006. Under uniform negatives, popularity alone already
reaches 0.731 on MovieLens, and on Yelp it ties TF-IDF (0.465) with every
encoder within 0.007 of it. This shows how much of an easy-negative
evaluation a non-personalized model can absorb \cite{krichene-sampled}.

\subsection{The capacity ladder on MovieLens}
\label{sec:cross-dataset-synthesis}
\begin{table}[t]
\centering
\caption{Paired per-user tests against SBERT, all datasets and scorers
(personalized scorer, popularity-matched negatives, 5 seeds). Cells are
NDCG@10 of the pipeline minus SBERT. $^{*}$: Holm-corrected $p<0.05$;
$^{\equiv}$: TOST-equivalent within $\pm0.01$; unmarked: inconclusive.}
\label{tab:scorer-comparison}
\resizebox{\columnwidth}{!}{\begin{tabular}{lrrrrrr}
\toprule
$-$ SBERT & \multicolumn{2}{c}{MovieLens} & \multicolumn{2}{c}{Amazon} & \multicolumn{2}{c}{Yelp} \\
\cmidrule(lr){2-3}\cmidrule(lr){4-5}\cmidrule(lr){6-7}
 & LGBM & MLP & LGBM & MLP & LGBM & MLP \\
\midrule
TF-IDF & -0.014$^{*}$ & -0.026$^{*}$ & -0.005 & -0.043$^{*}$ & -0.013$^{*}$ & -0.015$^{*}$ \\
mxbai (335M) & -0.007 & -0.008 & +0.003 & -0.015 & +0.003$^{\equiv}$ & +0.010$^{*}$ \\
phi4-mini raw (3.8B) & -0.104$^{*}$ & -0.135$^{*}$ & -0.015 & -0.101$^{*}$ & +0.002$^{\equiv}$ & -0.037$^{*}$ \\
Qwen3-Emb.\ (4B) & -0.002 & -0.012 & +0.010 & -0.002 & +0.006$^{*}$ & +0.016$^{*}$ \\
Hybrid & -0.010$^{*}$ & -0.023$^{*}$ & -0.003 & -0.066$^{*}$ & +0.007$^{*}$ & -0.020$^{*}$ \\
\bottomrule
\end{tabular}
}
\end{table}

Table~\ref{tab:scorer-comparison} reads the personalized,
popularity-matched configuration as a ladder of capacity steps. On
MovieLens (first two columns), moving
from TF-IDF to SBERT (22M parameters) is a significant gain
($+0.014$, $p<0.001$). It is the one step that pays: TF-IDF only ties
PureSVD, while SBERT beats it ($+0.017$, $p=0.011$). No step above SBERT
produces a detectable gain. mxbai (15$\times$ larger) and
Qwen3-Embedding-4B (180$\times$ larger) land within 0.007 of SBERT in
point estimate with CIs straddling zero. With 610 users, however, TOST
cannot establish equivalence at $\pm0.01$ either. We therefore claim
\emph{no detectable gain} above SBERT, not equivalence.

\textbf{A raw LLM is worse than no LLM.} Mean-pooled phi4-mini is the
outlier. It is $0.104$ below SBERT ($p<10^{-60}$), below TF-IDF, and
below PureSVD, which uses no text at all. Qwen3-Embedding-4B has
nearly the same size and the same quantization, and it recovers fully.
The deficit therefore comes from using an untrained chat model's hidden
states as an embedding, not from LLM capacity. Concatenating TF-IDF to
phi4-mini (Hybrid) recovers most of the loss but stays significantly
below SBERT alone ($-0.010$, $p=0.037$).

\textbf{Two scorers, one ordering.} The pointwise MLP reproduces this
ordering: SBERT is above TF-IDF ($+0.026$, $p<10^{-4}$), mxbai and
Qwen3 are not detectably different from SBERT, and phi4-mini ($-0.135$)
and Hybrid ($-0.023$) are significantly worse. The scorers disagree
about absolute levels, as expected for a listwise versus a pointwise
objective, but not about which representations help.

\subsection{Review text: Amazon and Yelp}
\label{sec:amazon-results}
\begin{table*}[t]
\centering
\caption{Amazon Reviews and Yelp NDCG@10 (LGBMRanker, mean $\pm$ 95\% CI
over 5 seeds). ``Item only'' and ``+Aff.+Pop.'' use popularity-matched
negatives, as in Table~\ref{tab:main-results}; ``Uniform'' uses uniform
negatives.}
\label{tab:amazon-results}
\begin{tabular}{lcccccc}
\toprule
 & \multicolumn{3}{c}{Amazon (249 users)} & \multicolumn{3}{c}{Yelp (10,537 users)} \\
\cmidrule(lr){2-4}\cmidrule(lr){5-7}
Pipeline & Item only & +Aff.+Pop. & Uniform & Item only & +Aff.+Pop. & Uniform \\
\midrule
Popularity & 0.090$\pm$0.003 & 0.090$\pm$0.003 & 0.132$\pm$0.012 & 0.133$\pm$0.001 & 0.133$\pm$0.001 & 0.465$\pm$0.003 \\
PureSVD & 0.151$\pm$0.007 & 0.151$\pm$0.007 & 0.166$\pm$0.007 & 0.151$\pm$0.003 & 0.151$\pm$0.003 & 0.298$\pm$0.002 \\
\midrule
TF-IDF & 0.068$\pm$0.009 & 0.170$\pm$0.011 & 0.180$\pm$0.005 & 0.110$\pm$0.003 & 0.179$\pm$0.001 & 0.465$\pm$0.003 \\
SBERT (22M) & 0.066$\pm$0.008 & 0.175$\pm$0.013 & 0.195$\pm$0.010 & 0.110$\pm$0.003 & 0.192$\pm$0.002 & 0.469$\pm$0.001 \\
mxbai (335M) & 0.064$\pm$0.009 & 0.178$\pm$0.015 & 0.187$\pm$0.010 & 0.109$\pm$0.002 & 0.195$\pm$0.001 & 0.472$\pm$0.002 \\
phi4-mini raw (3.8B) & 0.065$\pm$0.009 & 0.160$\pm$0.008 & 0.159$\pm$0.010 & 0.111$\pm$0.002 & 0.194$\pm$0.001 & 0.465$\pm$0.002 \\
Qwen3-Emb.\ (4B) & 0.065$\pm$0.010 & 0.185$\pm$0.007 & 0.196$\pm$0.013 & 0.110$\pm$0.002 & 0.198$\pm$0.002 & 0.469$\pm$0.003 \\
Hybrid & 0.066$\pm$0.009 & 0.172$\pm$0.008 & 0.186$\pm$0.007 & 0.110$\pm$0.002 & 0.199$\pm$0.001 & 0.470$\pm$0.003 \\
\bottomrule
\end{tabular}

\end{table*}

Table~\ref{tab:amazon-results} reports the two review datasets. Amazon
repeats the MovieLens pattern in direction but has little power. Every
personalized content pipeline beats popularity ($p<10^{-4}$), but with
249 evaluated users every LGBMRanker comparison between content
pipelines is inconclusive. The MLP scorer resolves only the extremes:
SBERT beats TF-IDF ($+0.043$) and phi4-mini is $0.101$ below SBERT.

Yelp, with 10{,}537 users, has enough power to resolve differences well
under 0.01. The first capacity step replicates: SBERT beats TF-IDF by
$0.013$ under LGBMRanker and $0.015$ under the MLP (both $p<10^{-12}$).
Unlike on the other datasets, capacity above SBERT also pays, but by
less. Qwen3-Embedding-4B is $+0.006$ above SBERT under LGBMRanker
($p=0.002$; 0.198 vs.\ 0.192, $+3\%$ relative) and $+0.016$ under the MLP
($+9\%$). mxbai is TOST-equivalent to SBERT under LGBMRanker and
$+0.010$ under the MLP. Raw phi4-mini is equivalent to SBERT under
LGBMRanker but $0.037$ below it under the MLP; \S\ref{sec:failure-analysis}
examines why. Across the three datasets, then, the TF-IDF-to-SBERT step
is the most consistent gain (on Amazon it is significant only under the
MLP). A trained LLM embedder adds a
further gain, half the size of that step under LGBMRanker and the same
size under the MLP, but only on long review text with a sample large
enough to detect it.

\subsection{Text redundancy, within one dataset}
\label{sec:redundancy-analysis}
\label{fig:redundancy-vs-gap}
Corpus lexical redundancy (mean pairwise TF-IDF cosine between items)
is a cheap candidate predictor of when dense encoders beat TF-IDF; in
retrieval, dense models' advantage over BM25 shrinks on high-overlap
corpora \cite{rare-redundancy}. Comparing datasets cannot isolate it,
because datasets also differ in domain, size, and text type, so we test
it within MovieLens. We
hold users, items, labels, splits, and negatives fixed and vary only
the item text:

\begin{center}
\small
\begin{tabular}{lccc}
\toprule
Item text & Redundancy & SBERT $-$ TF-IDF & $p$ \\
\midrule
Title only & 0.009 & $+0.019\pm0.008$ & $2\times10^{-6}$ \\
Title + genres & 0.030 & $+0.014\pm0.007$ & $5\times10^{-5}$ \\
Genres only & 0.188 & $-0.008\pm0.004$ & $3\times10^{-4}$ \\
\bottomrule
\end{tabular}
\end{center}

\noindent The dense encoder's advantage shrinks as redundancy rises and
reverses at the redundant extreme. All three differences survive a
Bonferroni correction. The manipulation is not a pure redundancy
intervention: genres-only text also removes the title, and with
$\sim$20 genre tokens TF-IDF becomes an exact genre indicator. The
result therefore shows that the lexical-versus-dense ordering
\emph{follows} text redundancy within a dataset, which is the property
a practitioner can measure before running anything. It does not show
that redundancy alone causes the ordering. Nor does it transfer across
datasets: Yelp's review text is the most redundant corpus we test
(0.236; Amazon 0.193), yet SBERT beats TF-IDF there by 0.013. Redundancy
is a guide for choosing which text to embed within one corpus, not a
predictor that ranks corpora.

\subsection{Robustness under temporal shift}
\label{sec:robustness}
\label{tab:robustness}
\label{fig:robustness}
Holding out each user's most recent positives lowers every pipeline
(popularity-matched negatives, LGBMRanker, 5 seeds). On MovieLens, NDCG@10
drops by 0.115--0.128, and the ordering established under the random
split holds. TF-IDF degrades most (0.576 $\rightarrow$ 0.448, a drop of 0.128),
while SBERT and mxbai degrade least (0.115 each; SBERT 0.589
$\rightarrow$ 0.474), so the SBERT-over-TF-IDF gap roughly doubles, from $+0.014$ to $+0.027$ ($p<10^{-4}$). mxbai also separates
from TF-IDF under shift ($+0.020$, $p=0.006$). No encoder above SBERT is
detectably better than it (mxbai $-0.006$, Qwen3 $-0.010$, both
$p>0.5$), and raw phi4-mini remains far behind ($-0.115$). The only
reliable robustness gain therefore again comes from the first step
off lexical features. On Amazon, every pipeline drops by a similar
amount (0.055--0.069), every content pipeline still beats popularity
($p<10^{-3}$), and no encoder is measurably more robust than another.
On Yelp the shift costs content pipelines only 0.005--0.009 (popularity
0.016). Both gains survive it: SBERT over TF-IDF ($+0.011$) and Qwen3
over SBERT ($+0.010$, $p<10^{-4}$). Among content pipelines, Qwen3 loses the least (0.005).

\subsection{Quality--cost}
\label{sec:quality-cost}

Table~\ref{tab:latency} reports per-item embedding latency from the
like-for-like benchmark (\S\ref{sec:experimental-setup}).

\begin{table}[t]
\centering
\caption{Median per-item embedding latency (ms) over three rounds, 256
items, full-text coverage. Ratios are relative to SBERT at the same
batch size. MovieLens text is a title and genres; Amazon text averages
about 2{,}100 characters of review snippets.}
\label{tab:latency}
\footnotesize
\setlength{\tabcolsep}{2pt}
\begin{tabular}{lrrrr}
\toprule
 & \multicolumn{2}{c}{MovieLens} & \multicolumn{2}{c}{Amazon} \\
\cmidrule(lr){2-3}\cmidrule(lr){4-5}
Encoder & batch 1 & batch 32 & batch 1 & batch 32 \\
\midrule
SBERT (22M) & 12.3 & 5.8 & 60 & 42 \\
mxbai (335M) & 42.5 (3.5$\times$) & 28.8 (5.0$\times$) & 257 (4.3$\times$) & 241 (5.7$\times$) \\
phi4-mini (3.8B) & 222 (18$\times$) & 170 (29$\times$) & 2{,}691 (45$\times$) & 2{,}696 (64$\times$) \\
Qwen3-Emb.\ (4B) & 231 (19$\times$) & 262 (45$\times$) & 2{,}974 (49$\times$) & 2{,}924 (70$\times$) \\
\bottomrule
\end{tabular}
\end{table}

With the same inputs, batching, and hardware for every encoder, the
two 4B-class models cost 18--45$\times$ SBERT per item on short
MovieLens text and 45--70$\times$ on long Amazon text, and mxbai costs
3.5--5.7$\times$. The gap widens with text length because SBERT and
mxbai process short chunks while the LLMs attend over the full text,
and widens with batch size because SBERT gains more from batching
than the LLM servers do. On Amazon, batching barely helps the LLMs:
with four parallel server slots on one machine, long sequences
saturate the accelerator even at batch size 1. Qwen3 at batch 32 on
MovieLens is slower per item than at batch 1, which we attribute to
contention among the server slots, a property of this deployment
rather than of the model. A naive comparison of SBERT's batched corpus
encoding against unbatched, sequential LLM calls inflates the ratio by
more than an order of magnitude, so latency claims should state the
batching regime.

Set against quality, the extra cost buys nothing measurable. On
MovieLens, Qwen3 costs 19--45$\times$ SBERT for an NDCG@10 difference of
$-0.002$ (inconclusive), mxbai costs 3.5--5$\times$ for $-0.006$, and
phi4-mini costs 18--29$\times$ for $-0.104$. On Amazon, Qwen3 has the
highest point estimate ($+0.010$ over SBERT under LGBMRanker, $-0.002$
under the MLP), but neither difference is detectable at 45--70$\times$
the cost. Yelp is the one case where capacity above SBERT buys a
detectable gain: $+0.006$ to $+0.016$ NDCG@10 for Qwen3. Yelp's item text
is as long as Amazon's (mean 2{,}148 vs.\ 2{,}083 characters), so its
per-item cost should be similar, about 45--70$\times$ SBERT. Because item embeddings are computed once per item and
cached, these ratios bound refresh and cold-start cost rather than
request latency (\S\ref{sec:practical-guidelines}).
